% Source: https://artofproblemsolving.com/community/c3h1764322_geometry
% Stable identity: aops-topic-1764322-q01
\textbf{Area of $ABC$.} By Heron's formula with $s=\dfrac{17+25+26}{2}=34$,
\[
[ABC]=\sqrt{34\cdot17\cdot9\cdot8}=\sqrt{41616}=204 .
\]

\textbf{The median triangle has $\dfrac{3}{4}$ of this area.} Let $G$ be the centroid and $D$ the midpoint of $BC$, and extend $GD$ beyond $D$ to $H$ with $DH=GD$.

\begin{asy}
usepackage("amsmath");
size(7cm);
defaultpen(fontsize(10pt));
pair B=(0,0),C=(26,0),A=(340/52,408/26);
pair D=(B+C)/2,G=(A+B+C)/3,H=2*D-G;
draw(A--B--C--cycle,linewidth(0.9));
draw(A--D,gray);draw(B--(A+C)/2,gray);draw(C--(A+B)/2,gray);
draw(G--H,blue+linewidth(0.8));draw(B--H,blue+dashed);draw(C--H,dashed+gray);
draw(B--G,blue+linewidth(0.8));
dot(G);dot(D);dot(H);
label("$A$",A,N);label("$B$",B,W);label("$C$",C,E);
label("$G$",G,dir(70));label("$D$",D,dir(-125));label("$H$",H,S);
label("$17$",A+0.3*(B-A),W);label("$25$",A+0.3*(C-A),NE);
label("\mbox{}",(min(currentpicture,user=true)+max(currentpicture,user=true))/2,black); // keeps the text after the figure black when LaTeX typesets the labels inline
\end{asy}

Since the diagonals $BC$ and $GH$ of $BGCH$ bisect each other, $BGCH$ is a parallelogram, so $BH=GC$. The centroid divides each median in the ratio $2:1$, so
\[
BG=\dfrac{2}{3}m_b,\qquad GH=2GD=\dfrac{2}{3}m_a,\qquad BH=GC=\dfrac{2}{3}m_c ,
\]
where $m_a,m_b,m_c$ are the medians from $A,B,C$. So triangle $BGH$ is similar to the median triangle with ratio $\dfrac{2}{3}$. Its area is $[BGH]=2[BGD]=[BGC]=\dfrac{1}{3}[ABC]$ (the median $BD$ of triangle $BGH$ halves it, and $D$ is the midpoint of $BC$). Hence the median triangle has area
\[
\left(\dfrac{3}{2}\right)^2\cdot\dfrac{1}{3}[ABC]=\dfrac{3}{4}\cdot204=\boxed{153}.
\]
